% ===================================================================== front matter
% ---------------------------------------------------------------- 1. cover page
\begin{titlepage}
\thispagestyle{empty}
\centering
\vspace*{8mm}
{\sffamily\large INTERNSHIP ENGINEERING REPORT\par}
\vspace{18mm}
{\sffamily\bfseries\Huge FLOW BEHAVIOR AND THERMAL EFFECTS IN MULTIPHYSICS SYSTEMS\par}
\vspace{10mm}
{\sffamily\Large CFD, Conjugate Heat Transfer and Thermo-Structural Analysis of a Heated Cylindrical Duct\par}
\vspace{16mm}
\includegraphics[width=0.78\linewidth]{PHYSICS_SCHEMATIC.png}\par
{\footnotesize\itshape Physics chain of the model (Section 3 drawing; values printed in the process boxes are Section 2 analytical estimates).\par}
\vspace{14mm}
{\large Rohan Balram Patel\par}
{\normalsize B.Tech Aerospace Engineering, Dayananda Sagar University, Bengaluru\par}
\vspace{6mm}
{\normalsize Eleation\par}
{\normalsize Internship: February--May 2025\par}
\vspace{10mm}
{\sffamily\small Numerical analysis using ANSYS Workbench, ANSYS Fluent and ANSYS Mechanical\par}
\vfill
\end{titlepage}
\pagenumbering{roman}

% ---------------------------------------------------------------- 2./3. internship information and project title
\chapter*{Internship Information}
\addcontentsline{toc}{chapter}{Internship Information}
\begin{tabular}{@{}>{\bfseries}p{0.30\linewidth}p{0.64\linewidth}@{}}
Engineer & Rohan Balram Patel \\
Programme & B.Tech Aerospace Engineering, Dayananda Sagar University, Bengaluru \\
Organisation & Eleation \\
Internship period & February 2025 -- May 2025 \\
Tools & ANSYS Workbench (SpaceClaim), ANSYS Fluent, ANSYS Mechanical / MAPDL; Python 3.10 (bundled with ANSYS) for calculations and post-processing \\
Report status & Final internship engineering report
\end{tabular}

\section*{Project Title}
\textbf{Flow Behavior and Thermal Effects in Multiphysics Systems}\\
CFD, Conjugate Heat Transfer and Thermo-Structural Analysis of a Heated Cylindrical Duct

% ---------------------------------------------------------------- 4. declaration / note on the analysis record
\chapter*{Declaration}
\addcontentsline{toc}{chapter}{Declaration}
\recbox{\textbf{Note on the analysis record.} The original internship project files were not retained. The geometry, operating
parameters, simulation set-up and numerical results presented in this report therefore come from a complete re-analysis of the
internship problem, carried out after the internship with ANSYS Workbench, Fluent and Mechanical and based on the documented project
scope. They are not copies of the original internship analysis.}

Only documented internship information (engineer, programme, organisation, period and tool set) is described as historical fact.
Every other value in this report is labelled with one of the following terms:

\begin{itemize}
\item \textbf{Selected} -- an input value chosen for this study with a stated rationale (for example $D_i$ = 20~mm, $V_{in}$ = 23.5~m/s, $q''$ = 8000~W/m$^2$, Inconel 718);
\item \textbf{Assumed} -- an engineering idealisation (for example steady state, uniform heat flux, $\nu$ = 0.294, idealised supports);
\item \textbf{Calculated} -- a hand or script calculation from the selected inputs (the Section 2 analytical baseline);
\item \textbf{Simulated} -- the output of an ANSYS solution of this project.
\end{itemize}

No experimental or measured data exist for this project, and none has been created. Verification is against analytical theory,
published correlations and internal consistency checks; the results are therefore verified but not experimentally validated.

\textbf{Declaration.} This report presents the analysis as documented in the project record (the file
\texttt{PROJECT\_STATE.md}, Sections 0--10A) and in the audited master dataset of Section 10A
(\texttt{MASTER\_PROJECT\_DATA.csv}). Every numerical result quoted is a
calculation or simulation of the analysed problem and is traceable to a project file. No original company procedure, supervisor
comment or original screenshot is reproduced or implied.

\vspace{6mm}
Rohan Balram Patel

% ---------------------------------------------------------------- 5. acknowledgement
\clearpage   % keeps the original page break (Acknowledgement and Abstract on page ii)
\chapter*{Acknowledgement}
\addcontentsline{toc}{chapter}{Acknowledgement}
I thank Eleation for the internship opportunity between February and May 2025, during which the project on flow behaviour and
thermal effects in multiphysics systems was carried out, and the Department of Aerospace Engineering of Dayananda Sagar University,
Bengaluru, for the academic framework of the internship. The analysis documented here relied on ANSYS Workbench, Fluent and Mechanical
and on the published data and correlations listed in the references.

% ---------------------------------------------------------------- 6. abstract
\chapter*{Abstract}
\addcontentsline{toc}{chapter}{Abstract}
This report presents the Eleation internship project \emph{Flow Behavior and Thermal Effects in Multiphysics Systems}
(February--May 2025): a numerical analysis using ANSYS Workbench, ANSYS Fluent and ANSYS Mechanical. The system is an air-cooled,
externally heated Inconel 718 duct (20/40~mm inner/outer diameter, 600~mm long).

An analytical baseline preceded all simulations. Turbulent internal flow and conjugate heat transfer were solved in ANSYS Fluent (steady,
$k$--$\omega$ SST, wall-resolved) at 23.5~m/s and 300~K inlet with an outer-wall heat flux of 8000~W/m$^2$. On the 159,840-cell
medium mesh the converged solution gives a pressure drop of 438.13~Pa, an outlet bulk temperature of 368.93~K and a maximum
solid temperature of 562.58~K, with only 7.58~K across the wall at mid-span. A three-mesh study shows approximately convergent behaviour
(fine mesh 560.83~K; extrapolated 558.13~K).

The solid temperature field was transferred one-way to ANSYS Mechanical (mapping error $\le$ 0.081~K). Free expansion gives 1.84~mm
growth and a 24.28~MPa peak von Mises stress. Full axial restraint gives a 548.94~kN end force, a mean axial stress of $-$582.44~MPa and
a local peak of 605.16~MPa, 57.8\,\% of the local yield strength; structural mesh effects are below
$5\times10^{-4}$. Linear eigenvalue buckling under the baseline S1 support gives $\lambda_1$ = 1.108, below the first-yield factor of 1.730,
while alternative idealised supports give 2.232 (S3) and 4.300 (S2). A one-factor parametric study (velocity and heat flux $\pm$10\,\%,
wall thickness 8 and 12~mm), with full solutions of every case, shows that velocity acts mainly on cooling and pressure drop,
heat flux drives temperature and stress, and thickness drives stiffness and buckling capacity; $\lambda_1$ falls below 1 for +10\,\% heat
flux and the 8~mm wall. An additional geometrically nonlinear analysis with ANSYS LS-DYNA, performed as an extension with mode-shaped
numerical imperfections of 0.1, 0.6 and 1.2~mm, gives Southwell characteristic-load estimates of 607.7--612.3~kN, close to the linear
critical load of 608.25~kN, while lateral deflection and stress grow strongly with the imperfection amplitude.

The results are not validated against measurement. The real end restraint, measured geometric imperfections and
inelastic behaviour, the dominant open uncertainties, were not modelled. The study therefore characterises the thermal-stress mechanism and its sensitivities, not the
structural adequacy of a real component.

% ---------------------------------------------------------------- 7.-9. contents and lists
\cleardoublepage
\chapter*{Contents}
\addcontentsline{toc}{chapter}{Contents}
{\small\makeatletter\setlength{\columnsep}{8mm}\begin{multicols}{2}\@starttoc{toc}\end{multicols}\makeatother}
\clearpage
\phantomsection
\addcontentsline{toc}{chapter}{List of Figures}
{\sffamily\bfseries\LARGE List of Figures\par}\vspace{2pt}\hrule\vspace{8pt}
{\makeatletter\let\origaddvspace\addvspace\renewcommand{\addvspace}[1]{\origaddvspace{2pt}}\@starttoc{lof}\makeatother}
\vspace{14pt}
\phantomsection
\addcontentsline{toc}{chapter}{List of Tables}
{\sffamily\bfseries\LARGE List of Tables\par}\vspace{2pt}\hrule\vspace{8pt}
{\makeatletter\let\origaddvspace\addvspace\renewcommand{\addvspace}[1]{\origaddvspace{2pt}}\@starttoc{lot}\makeatother}

% ---------------------------------------------------------------- 10. abbreviations
\vspace{10pt}
\noindent\begin{minipage}{\linewidth}
\phantomsection
\addcontentsline{toc}{chapter}{List of Abbreviations and Symbols}
{\sffamily\bfseries\LARGE List of Abbreviations and Symbols\par}\vspace{2pt}\hrule\vspace{6pt}
{\footnotesize
\begin{tabular}{@{}p{0.10\linewidth}p{0.32\linewidth}p{0.12\linewidth}p{0.38\linewidth}@{}}
\toprule
\multicolumn{2}{@{}l}{\textbf{Abbreviations}} & \multicolumn{2}{l}{\textbf{Symbols}} \\
\midrule
CAD & computer-aided design & $A$, $I$ & section area [m$^2$]; second moment of area [m$^4$] \\
CFD & computational fluid dynamics & $c_p$ & specific heat [J/(kg\,K)] \\
CHT & conjugate heat transfer & $D_i$, $D_o$, $D_h$ & inner, outer, hydraulic diameter [m] \\
FE, FV & finite element, finite volume & $E$, $\nu$ & Young's modulus [Pa]; Poisson's ratio [--] \\
GCI & grid convergence index & $f$ & Darcy friction factor [--] \\
LC1 & load case 1, free expansion & $h$ & heat-transfer coefficient [W/(m$^2$\,K)] \\
LC2 & load case 2, axially restrained (S1) & $k$, $k_s$ & conductivity of air, of the solid [W/(m\,K)] \\
LC2P & LC2 with bore pressure & $K$, $L$ & effective-length factor [--]; length [m] \\
MAPDL & Mechanical APDL solver & $\dot m$, $N$ & mass flow rate [kg/s]; axial end force [N] \\
OQ & orthogonal quality & $Nu$, $Pr$, $Re$ & Nusselt, Prandtl, Reynolds numbers \\
P00 & official baseline case & $P_{cr}$ & critical buckling load [N] \\
RANS & Reynolds-averaged Navier--Stokes & $q''$, $Q$ & heat flux [W/m$^2$]; heat-transfer rate [W] \\
SST & shear-stress transport & $r$, $t$ & radius, radius of gyration [m]; wall thickness [m] \\
S1--S3 & support scenarios (Table~\ref{tab:supports}) & $S_y$ & 0.2\,\% proof (yield) strength [Pa] \\
TI & turbulence intensity & $T$, $T_{ref}$ & temperature, stress-free reference [K] \\
VM & von Mises (equivalent) stress & $y^+$ & dimensionless wall distance [--] \\
V01, V03 & velocity cases & $\alpha$ & thermal expansion coefficient [1/K] \\
Q01, Q03 & heat-flux cases & $\lambda_1$ & first linear buckling load factor [--] \\
T01, T03 & wall-thickness cases & $\mu$, $\rho$ & dynamic viscosity [Pa\,s]; density [kg/m$^3$] \\
1-D & one-dimensional & $\sigma_{vM}$, $\bar\sigma_z$ & von Mises stress; mean axial stress [Pa] \\
\bottomrule
\end{tabular}}
\end{minipage}
